\section{Trabalhos Relacionados}

Foram selecionados seis estudos empíricos de Engenharia de Software que avaliam PBT, a relação entre cobertura e \MT{} ou ferramentas de mutação; quatro deles foram publicados entre 2023 e 2025. Para cada um, apresentam-se objetivo e método, os resultados relevantes e a relação com este trabalho. A Tabela~\ref{tab:relacionados} resume o posicionamento.

\textbf{\citet{ravi2025pbt}} avaliaram a efetividade de PBT em Python. Analisaram 426 programas que usam a biblioteca Hypothesis, definiram 12 categorias de propriedade e aplicaram \MT{} às suítes de 40 projetos. Cada teste de propriedade encontrou, em média, cerca de 50 vezes mais mutações que um teste unitário; testes que verificam exceções, pertinência e tipos foram os mais efetivos; e 76\% das mutações encontradas apareceram nas primeiras 20 entradas geradas. É o trabalho mais próximo deste. Aqui, a mesma pergunta é feita em dois ecossistemas, com propriedades controladas, e o efeito do número de entradas é medido diretamente (10, 100 e 1000 tentativas).

\textbf{\citet{goldstein2024pbt}} investigaram como o PBT é usado na prática por meio de 30 entrevistas semiestruturadas na Jane Street, empresa que usa PBT em OCaml, com análise temática. Os participantes concentraram o uso em poucos padrões de alto retorno, como comparação com um modelo de referência e propriedades de ida e volta, e relataram dificuldade em escrever geradores e em avaliar a efetividade das suítes. Este trabalho responde a essa dificuldade de forma quantitativa: mede, por mutação, quanto propriedades com modelo de referência detectam a mais que propriedades de limite.

\textbf{\citet{corgozinho2023pbt}} caracterizaram como desenvolvedores implementam propriedades, classificando 86 testes de propriedade de 30 repositórios Python. Propriedades de ida e volta (37,2\%) e de oráculo (33,7\%) somaram 70,9\% dos testes. As categorias catalogadas por esse estudo orientaram a escolha das propriedades da suíte fraca deste trabalho, cuja efetividade é agora medida.

\textbf{\citet{jain2023gap}} propuseram a ``lacuna de oráculo'', a diferença entre cobertura e \ms{}, e a estudaram em 45 projetos Java, em arquivos com asserções removidas e em código de criptomoedas. Sobre o código coberto, cobertura e \ms{} tiveram correlação fraca (0,21), e remover asserções ampliou a lacuna. A QP5 deste trabalho mede essa lacuna em PBT e mostra como ela se fecha quando o oráculo das propriedades fica mais forte.

\textbf{\citet{kintis2018tools}} compararam quatro ferramentas de mutação para Java (PIT, muJava, Major e PITRV) com análise manual dos mutantes e com defeitos reais. A efetividade variou muito entre as ferramentas, e nenhuma subsumiu as demais. \textbf{\citet{monteiro2022tools}} compararam PIT e Major: o PIT gerou significativamente mais mutantes, com desempenho um pouco melhor, mas o Major superou o PIT quando os mutantes possivelmente equivalentes foram excluídos. Esses resultados motivam o recorte harmonizado de operadores deste trabalho, que estende a comparação de ferramentas ao par PIT/StrykerJS, entre linguagens.

\begin{table}[!htbp]
\centering
\caption{Posicionamento dos trabalhos relacionados.}
\label{tab:relacionados}
\small
\begin{tabularx}{\textwidth}{@{}L{3.4cm} Y c c c@{}}
\toprule
Trabalho & Método e linguagem & PBT & Mutação & Compara ecossistemas \\
\midrule
\citet{ravi2025pbt} & Mineração e mutação, Python & sim & sim & não \\
\citet{goldstein2024pbt} & Entrevistas, OCaml & sim & não & não \\
\citet{corgozinho2023pbt} & Mineração, Python & sim & não & não \\
\citet{jain2023gap} & Cobertura $\times$ mutação, Java & não & sim & não \\
\citet{kintis2018tools}; \citet{monteiro2022tools} & Comparação de ferramentas, Java & não & sim & não \\
\textbf{Este trabalho} & Experimento, Java e Node.js & sim & sim & sim \\
\bottomrule
\end{tabularx}
\end{table}
